% Source: https://artofproblemsolving.com/community/c3h2506496_is_this_true
% Stable identity: aops-topic-2506496-q01
\textbf{(a)} \textbf{No.} Take $a_1=3$, $b_1=1$, $c_1=2$ (so $3\equiv1\pmod2$) and $a_2=b_2=1$, $c_2=3$. Then $a_1a_2=3$ and $b_1b_2=1$, but $3\not\equiv1\pmod6$. (Even with $c_1=c_2$ it fails: $4\equiv1\pmod3$ and $1\equiv1\pmod3$, but $4\not\equiv1\pmod9$.)

\textbf{(b)} \textbf{No.} Take $c_1=c_2=2$, $a_1=1$, $b_1=0$ and $a_2=b_2=0$. Then $a_1a_2=0=b_1b_2$, so certainly $a_1a_2\equiv b_1b_2\pmod4$, but $1\not\equiv0\pmod2$.

\textbf{(c)} Write
\[a_1a_2-b_1b_2=a_1(a_2-b_2)+b_2(a_1-b_1).\]
Since $d\mid c_2\mid a_2-b_2$ and $d\mid c_1\mid a_1-b_1$, both terms are divisible by $d$, so $a_1a_2\equiv b_1b_2\pmod d$. In particular, if $c_1=c_2=c$, congruences mod $c$ can be multiplied. Finally, $a_1-b_1=c_1k$ and $a_2-b_2=c_2l$ for integers $k,l$, so $(a_1-b_1)(a_2-b_2)=c_1c_2kl$. $\blacksquare$

So what can be multiplied modulo $c_1c_2$ is the product of the \emph{differences}, not the difference of the products.
